\begin{solapartat}
\punts{0,5} Com $\lambda = 0,136\ \mathrm{m}$ i $v = 340\ \mathrm{m/s}$, el vector d'ones, la freqüència i la freqüència angular s'expressen com:
\[ k = \frac{2\pi}{\lambda} = 46,2\ \mathrm{m}^{-1} \]
\[ f = \frac{1}{T} = \frac{v}{\lambda} = 2500\ \mathrm{Hz} \]
\[ \omega = 2\pi f = 15\,710\ \mathrm{rad/s} \]

\punts{0,75} A partir de les relacions anteriors, podem expressar una ona harmònica que es propaga cap endavant de diferents maneres equivalents:
\[ \psi(x,t) = I_0\sin k(x-vt) \]
\[ \psi(x,t) = I_0\sin 2\pi\left(\frac{x}{\lambda}-\frac{t}{T}\right) \]
\[ \psi(x,t) = I_0\sin(kx-\omega t) \]
\[ \psi(x,t) = I_0\sin 2\pi f\left(\frac{x}{v}-t\right) \]

Finalment obtenim:
\[ A(x,t) = A_0\sin(46,2\,x-15\,710t),\ x\ \text{en m i}\ t\ \text{en s.} \]
\end{solapartat}

\begin{solapartat}
\punts{0,25} percep un so més agut, un augment de la freqüència que està associat a l'efecte Doppler.

\punts{0,2} L'explicació correcta és la segona: ``L'aparent canvi de la freqüència és degut al moviment relatiu entre la font i l'observador''

\punts{0,2} Es descarta la 1 perquè la ona sonora no canvia de freqüència, es tracta d'una apreciació i a més la font està en repòs.

\punts{0,2} Es descarta la 3 perquè no canvia de longitud d'ona de la ona sonora.

\punts{0,2} Es descarta la 4 perquè no canvia de longitud d'ona de la ona sonora. A més, el que és rellevant és el moviment relatiu entre l'observador i la font.

\punts{0,2} L'explicació 2 és la que descriu millor el fenomen. L'observador aprecia un canvi en la freqüència (aparent canvi en la freqüència) perquè el que distingeix l'oïda és la freqüència dels sons, com s'acosta a la font la separació temporal entre les crestes de la ona disminueix provocant aquest augment aparent de la freqüència. Per altra banda defineix molt bé el moviment, moviment relatiu; no determina que ni la font ni l'observador estan en repòs.
\end{solapartat}
